\section{Question}

This work began with a practical frustration. I was building an agent that wrote programs in a domain-specific language and repaired them using compiler feedback. The feedback arrived as text, and the model routinely failed to act on it: it would acknowledge an error and then resubmit the same code, or fall back to what its priors suggested, or simply give up. I first tried managing the agent's context so that the important feedback stayed at high priority, but the model still ignored it often enough to be unreliable. That failure changed the question I was asking. Instead of asking how to phrase feedback so that the model listens, I asked what happens if listening is no longer optional, that is, if the guarantee is moved out of the prompt and into the infrastructure that delivers the feedback.

The experiment compares two agent designs on a simple repair task. The model is given a short Python function into which exactly one bug has been introduced by changing a single line, and its job is to correct the code. The function is presented bare: its description, docstrings, and comments have been stripped beforehand, so the model cannot read what the code is supposed to do. The only way for the model to recover the missing specification is to submit code and learn from the checker's feedback. The two designs are identical in every respect except one. In both, the model submits code, a checker executes the task's tests, and the verdict is returned to the model as feedback, with the feedback text byte-identical across the two designs. The difference is who holds the authority to end the episode. In advisory mode, the model may declare the task done, and the declaration is believed even when the last verdict was FAILED. In binding mode, the declaration is ignored as mere text, and the episode ends as solved only when the checker passes. In other words, completion is computed by the checker rather than claimed by the model.

\begin{figure}[t]
\centering
\includegraphics[width=\textwidth]{fig7_modes.png}
\caption{The two modes: the same loop between the LLM and the checker, with the authority to end it held by the LLM in advisory and by the checker in binding. Drawn by \texttt{writeup\_figures/plot\_diagrams.py}.}
\label{fig:modes}
\end{figure}

The failure that separates the two designs is the false-DONE, in which the model submits wrong code and, in the same turn, confidently declares the task finished without waiting for the verdict. This is neither laziness nor defiance of feedback, because the mistake happens earlier than the feedback. Even though the checker is present and available, the model does not wait to consult it. To judge its own work, the model needs the specification that was stripped away, so it fills the gap from its prior knowledge: it has seen thousands of similar functions, so it imposes their learned conventions and constructs the missing specification for itself. It then grades its work against that constructed answer key, finds everything satisfied, and declares done. The answer key is wrong, and the model has no way to know it is wrong, because its only source of task knowledge produced the wrong task. The model does not know when it does not know.

This matters beyond the laboratory because in real deployments the specification is always incomplete. No user writes down every edge case and every convention. An agent is therefore always partly grading itself against a reconstructed answer key, and the question of whether its own ``done'' should carry authority is a design decision every agent system makes, usually without noticing.
